\documentclass[12pt]{article}
\usepackage[utf8]{inputenc}
\usepackage[T2A]{fontenc}
\usepackage[russian]{babel}
\usepackage[a4paper,margin=2cm]{geometry}
\usepackage{xcolor}
\usepackage{hyperref}
\usepackage{enumitem}
\usepackage{array}
\hypersetup{
    colorlinks=true,
    linkcolor=blue,
    filecolor=blue,
    urlcolor=blue,
    citecolor=blue,
}

\begin{document}

\begin{center}
    {\Large \textbf{Договор N 1 \\ купли-продажи программного обеспечения}} \\[0.5em]
    Версия документа: 1.0 \\
    Дата: 11 Июля, 2025 г.
\end{center}

\noindent
Ильичев Павел Юрьевич, именуем в дальнейшем ``Покупатель'', с одной стороны, и Меньшиков Никита Михайлович, именуем в дальнейшем ``Продавец'', с другой стороны, вместе именуемые в дальнейшем ``Стороны'', заключили настоящий Договор о нижеследующем.

\section{Предмет договора}

\begin{enumerate}
    \item Продавец обязуется передать в собственность Покупателю программный продукт — «Сервис определения маркировки и размера шины по фото» (далее — «Программный продукт»), разработанный в соответствии с Техническим заданием (Приложение №1 к настоящему Договору), а Покупатель обязуется принять Программный продукт и уплатить за него цену, установленную настоящим Договором.

    \item Наименование, характеристика и количество Программного продукта:
          \begin{enumerate}
              \item Наименование: «Сервис определения маркировки и размера шины по фото».
              \item Характеристика: микросервис, реализующий функции, описанные в разделах о сервисе ``Определение маркировки и размера шины по фото'' Технического задания, включая приём изображения шины и выдачу марки, модели и размерности шины в формате JSON; развёртывание в Docker-контейнере, взаимодействие через REST API с аутентификацией по API-ключу.
              \item Количество: 1 (одна) лицензия на право использования Программного продукта без ограничения по сроку и территории.
          \end{enumerate}

    \item Право собственности на Программный продукт переходит к Покупателю с момента подписания Сторонами Акта приёмки-передачи Программного продукта.
\end{enumerate}

\section{Обязанности сторон}

\begin{enumerate}
    \item Продавец обязан:
          \begin{enumerate}
              \item Передать Покупателю Программный продукт в течение 5 (пяти) рабочих дней с момента подписания Акта приёмки-передачи.
              \item Передать вместе с Программным продуктом следующую документацию: исходный код, Docker-файл(ы), файл \texttt{requirements.txt}, пользовательскую и техническую документацию (описание API, инструкция по развёртыванию), соответствующие Техническому заданию.
              \item Обеспечить техническую поддержку Программного продукта в течение 6 (шести) месяцев с даты ввода в эксплуатацию, в объёме, определённом разделом «Гарантийная поддержка» Технического задания.
              \item Не разглашать данные, полученные в ходе работы над проектом, а также не предоставлять третьим лицам доступ к исходному коду и весам моделей без письменного согласия Покупателя.
          \end{enumerate}

    \item Покупатель обязуется:
          \begin{enumerate}
              \item Принять Программный продукт по Акту приёмки-передачи.
              \item Оплатить Программный продукт в сроки, установленные Разделом 3 настоящего Договора.
              \item Предоставить Продавцу необходимые данные и доступы, требуемые для ввода Программного продукта в эксплуатацию.
          \end{enumerate}
\end{enumerate}

\section{Финансовые условия и порядок расчётов}

\begin{enumerate}
    \item Стоимость Программного продукта составляет 250 000 (двести пятьдесят тысяч) рублей.
    \item Порядок оплаты:
          \begin{enumerate}
              \item 50 000 (пятьдесят тысяч) рублей — в день подписания настоящего Договора.
              \item 50 000 (пятьдесят тысяч) рублей — не позднее 14 (четырнадцати) календарных дней с момента подписания настоящего Договора.
              \item 50 000 (пятьдесят тысяч) рублей — не позднее 7 (семи) календарных дней с момента фактического ввода Программного продукта в эксплуатацию.
              \item 100 000 (сто тысяч) рублей — не позднее 30 ноября 2025 года.
          \end{enumerate}
    \item Оплата производится на расчётный счёт Продавца, указанный в реквизитах, безналичным переводом.
\end{enumerate}

\section{Гарантийный срок программного продукта}

\begin{enumerate}
    \item Гарантийный срок Программного продукта составляет 6 (шесть) месяцев.
    \item Гарантийный срок начинается с даты подписания Акта приёмки-передачи Программного продукта.
    \item Если Покупатель лишён возможности использовать Программный продукт по обстоятельствам, зависящим от Продавца, гарантийный срок приостанавливается на период действия указанных обстоятельств.
    \item Гарантийный срок продлевается на время, в течение которого Программный продукт не мог использоваться из-за обнаруженных в нём недостатков, при условии извещения Продавца о недостатках в течение 5 (пяти) рабочих дней с момента их обнаружения.
\end{enumerate}

\section{Ответственность сторон}

\begin{enumerate}
    \item Имущественная и иная ответственность Сторон определяется настоящим Договором и действующим законодательством РФ.
    \item В случае нарушения Покупателем срока платежа Продавец вправе потребовать уплаты неустойки в размере 0,1\% от суммы просроченного платежа за каждый день просрочки.
    \item Выплата неустоек не освобождает Стороны от исполнения своих обязательств.
\end{enumerate}

\section{Обстоятельства непреодолимой силы (форс-мажор)}

\begin{enumerate}
    \item Порядок освобождения Сторон от ответственности при наступлении обстоятельств непреодолимой силы определяется действующим законодательством РФ. Сторона, подвергшаяся воздействию таких обстоятельств, обязана уведомить другую Сторону в письменной форме в течение 5 (пяти) календарных дней.

    \item Сторона, которая подвергается воздействию непреодолимой силы, должна доказать существование непреодолимой силы достоверными документами.
\end{enumerate}

\section{Порядок разрешения споров}

Все споры и разногласия, возникающие из настоящего Договора, решаются путём переговоров. В случае недостижения согласия спор подлежит рассмотрению в суде в соответствии с законодательством РФ.

\section{Срок действия договора}

\begin{enumerate}
    \item Настоящий Договор вступает в силу с момента подписания и действует до полного исполнения Сторонами своих обязательств.
    \item Прекращение действия Договора не освобождает Стороны от ответственности за нарушения, имевшие место в период его действия.
\end{enumerate}

\section{Заключительные положения}

\begin{enumerate}
    \item Все приложения и дополнительные соглашения к настоящему Договору являются его неотъемлемой частью и действительны при условии их письменного оформления и подписания уполномоченными представителями Сторон.
    \item Стороны обязаны незамедлительно письменно уведомлять друг друга об изменениях реквизитов.
    \item Во всём, что не предусмотрено настоящим Договором, Стороны руководствуются действующим законодательством РФ.
    \item Настоящий Договор составлен и подписан в двух экземплярах, имеющих равную юридическую силу — по одному для каждой из Сторон.
\end{enumerate}

\section{Реквизиты сторон}
\begin{center}
    \begin{tabular}{|>{\raggedright\arraybackslash}p{4cm}|>{\raggedright\arraybackslash}p{5cm}|>{\raggedright\arraybackslash}p{5cm}|}
        \hline
        \textbf{Реквизит}      & \textbf{Покупатель}      & \textbf{Продавец}           \\ \hline
        Имя                    & Ильичев Павел Юрьевич    & Меньшиков Никита Михайлович \\ \hline
        ИНН                    & 732896478438             & 524709285864                \\ \hline
        ОГРНИП                 & 322730000013731          &                             \\ \hline
        КПП                    &                          & 526002001                   \\ \hline
        Расчётный счёт         & 40802 810 0 0000 3397713 & 40817 810 9 4200 7326135    \\ \hline
        Банк                   & АО «ТБанк»               & ПАО «Сбербанк»              \\ \hline
        БИК                    & 044525974                & 042202603                   \\ \hline
        Корреспондентский счёт & 30101810145250000974     & 30101 810 9 0000 0000603    \\ \hline
    \end{tabular}
\end{center}

\vspace{1em}

\begin{flushleft}
    М.П. \hfill М.П.
\end{flushleft}

\end{document}
